\documentclass[11pt]{article}
\usepackage[margin=0.75in]{geometry}
\usepackage{fancyhdr}
\usepackage{array}
\usepackage{booktabs}
\usepackage{xcolor}
\usepackage{enumitem}

\pagestyle{fancy}
\fancyhf{}
\rhead{Patient: Santos, Maria E.}
\lhead{Riverside Imaging Center}
\rfoot{Page \thepage}

\definecolor{imagingpurple}{RGB}{75,0,130}

\begin{document}

\begin{center}
{\Large\bfseries\color{imagingpurple} RIVERSIDE IMAGING CENTER}\\[0.3em]
{\small 600 Diagnostic Drive, Riverside, CA 92501}\\
{\small Phone: (951) 555-4500 | Fax: (951) 555-4501}\\[1em]
{\LARGE\bfseries RADIOLOGY REPORT}
\end{center}

\vspace{0.5em}
\hrule height 2pt
\vspace{1em}

\section*{PATIENT INFORMATION}

\begin{tabular}{@{}p{3.5cm}p{5cm}p{3cm}p{4cm}@{}}
\textbf{Patient:} & Maria Elena Santos & \textbf{DOB:} & 03/14/1988 \\
\textbf{Exam Date:} & February 15, 2026 & \textbf{Accession:} & RIC-2026-021547 \\
\textbf{Ordering MD:} & Kevin Park, MD & \textbf{Insurance:} & Hartford WC \\
\textbf{Radiologist:} & Amanda Chen, MD & & \\
\end{tabular}

\vspace{1em}
\hrule
\vspace{1em}

\begin{center}
{\Large\bfseries MRI LUMBAR SPINE WITHOUT CONTRAST}
\end{center}

\section*{CLINICAL HISTORY}

35-year-old female with lower back pain following fall from scaffolding (15 feet) on 02/08/2026. Rule out disc herniation, fracture.

\section*{TECHNIQUE}

MRI of the lumbar spine was performed on a 1.5T magnet without intravenous contrast. Sagittal T1, T2, and STIR sequences were obtained along with axial T2-weighted images through the disc levels.

\section*{COMPARISON}

X-ray lumbar spine from 02/08/2026.

\section*{FINDINGS}

\textbf{Vertebral Bodies:}
\begin{itemize}[nosep]
\item Normal vertebral body height and alignment
\item No acute compression fracture identified
\item No marrow signal abnormality to suggest contusion or edema
\item Mild degenerative endplate changes at L4-L5
\end{itemize}

\textbf{Intervertebral Discs:}

\textit{L1-L2:} Normal disc height and signal. No herniation.

\textit{L2-L3:} Normal disc height and signal. No herniation.

\textit{L3-L4:} Mild disc desiccation. Small central disc bulge without significant canal stenosis. No nerve root impingement.

\textit{L4-L5:} \textbf{Moderate disc desiccation with posterior disc herniation.} There is a broad-based central and left paracentral disc protrusion measuring approximately 5mm, extending into the left lateral recess. This results in \textbf{moderate left lateral recess stenosis} with \textbf{contact of the traversing left L5 nerve root}. The central canal is mildly narrowed but not critically stenotic.

\textit{L5-S1:} Mild disc bulge without significant stenosis. No nerve root impingement.

\textbf{Conus Medullaris:} Normal in position, terminating at L1. No signal abnormality.

\textbf{Paraspinal Soft Tissues:} Mild edema within the left paraspinal musculature at L4-L5 level, consistent with muscle strain. No hematoma identified.

\textbf{Facet Joints:} Mild bilateral facet arthropathy at L4-L5 and L5-S1. No significant facet effusion.

\section*{IMPRESSION}

\begin{enumerate}
\item \textbf{L4-L5 posterior disc herniation with broad-based central and left paracentral protrusion}, resulting in moderate left lateral recess stenosis and contact with the left L5 nerve root. This is \textbf{likely acute or acutely worsened} given clinical history of recent trauma.

\item Mild paraspinal muscle strain at L4-L5 level.

\item No acute vertebral fracture.

\item Mild degenerative changes at L3-L4 and L5-S1 as described, likely pre-existing.
\end{enumerate}

\section*{RECOMMENDATION}

Clinical correlation recommended. If radicular symptoms persist, neurosurgical or pain management consultation may be warranted.

\vspace{2em}

\begin{tabular}{@{}p{10cm}@{}}
\hrulefill \\
\textbf{Amanda Chen, MD} \\
Board Certified Diagnostic Radiology \\
Riverside Imaging Center \\
Electronically signed: February 15, 2026 at 3:42 PM \\
\end{tabular}

\newpage

\begin{center}
{\Large\bfseries CT HEAD WITHOUT CONTRAST}\\
{\normalsize (From Emergency Department - February 8, 2026)}
\end{center}

\vspace{1em}

\section*{CLINICAL HISTORY}

Fall from height with head injury. Evaluate for intracranial hemorrhage.

\section*{TECHNIQUE}

Non-contrast CT of the head was performed with 5mm axial sections and coronal/sagittal reformats.

\section*{FINDINGS}

\textbf{Brain Parenchyma:}
\begin{itemize}[nosep]
\item No acute intracranial hemorrhage
\item No mass effect or midline shift
\item Gray-white matter differentiation preserved
\item No evidence of cerebral edema
\item Ventricles normal in size and configuration
\end{itemize}

\textbf{Extra-axial Spaces:}
\begin{itemize}[nosep]
\item No subdural or epidural hematoma
\item No subarachnoid hemorrhage
\end{itemize}

\textbf{Calvarium:}
\begin{itemize}[nosep]
\item No skull fracture identified
\item Soft tissue swelling noted overlying right frontal bone, consistent with contusion
\end{itemize}

\textbf{Visualized Paranasal Sinuses:} Clear.

\section*{IMPRESSION}

\begin{enumerate}
\item \textbf{No acute intracranial abnormality.}
\item \textbf{No skull fracture.}
\item Right frontal soft tissue swelling consistent with external trauma.
\item Clinical diagnosis of concussion supported; CT negative for structural injury requiring intervention.
\end{enumerate}

\vspace{2em}

\begin{tabular}{@{}p{10cm}@{}}
\hrulefill \\
\textbf{Michael Torres, MD} \\
Emergency Radiology \\
Riverside General Hospital \\
Electronically signed: February 8, 2026 at 11:58 AM \\
\end{tabular}

\end{document}
